% Beamer's default font size is 11 points. It is possible to set the
% default font size to any of 8, 9, 10, 11, 12, 14, 17, 20
\documentclass[13pt]{beamer}

\usetheme{Warsaw}
%\usecolortheme{dolphin}
\usecolortheme{seahorse}
%\useinnertheme{circles}
\usefonttheme{professionalfonts}
%\usefonttheme{structurebold}
%\usefonttheme{default}

\setbeamertemplate{navigation symbols}{}
 
\usepackage[english]{babel}
\usepackage[latin1]{inputenc}
\usepackage{times}
\usepackage[T1]{fontenc}
\usepackage{fancybox}

% These things can change:
\newcommand{\baseurl}{http://cell.vtt.fi/latex}
\urlstyle{sf}
\newcommand{\plus}{\makebox[2em]{+}}
\newcommand{\minus}{\makebox[2em]{-}}
\newcommand{\gray}[1]{%
\textcolor{gray}{#1}
}

% In beamer, defining new commands inside the slides fails.
% Definitions here:

\begin{document}

\logo{\includegraphics[width=1.4cm]{img/vttplain}\quad}
\title{\scalebox{1.5}{\textrm{\LaTeX}} Course 2011}
\subtitle{Part 3: Typographical commands}
\author{Arho Virkki}
\institute{\textsc{VTT Technical Research Centre of Finland}}
\date{\scalebox{0.8}{\textcolor{gray}{}}}

\begin{frame}
  \titlepage
\end{frame}


\begin{frame}[fragile]\frametitle{Elements}
The non-textual elements, including
\begin{itemize}
\item boxes (\verb|\makebox, \framebox ...|)
\item lines (\verb|\rule |)
\item distances (\verb|\hspace, \hspace*, \vspace, ...|)
\item \ldots
\end{itemize}
are described with commands inside the document.\medskip

Their arguments are given in millimeters, centimeters, pixels, or in proportion
of the font 'M' or 'x' characters,
\begin{center}
10mm, 1.5cm, 12px, 1em, 0.8ex, \ldots
\end{center}
\end{frame}


\begin{frame}[fragile]\frametitle{Examples}
Let us inspect some of the most common commands:

\begin{itemize}
\item \verb|\rule[voffset]{width}{height}| \\ 
  e.g.:  \verb|Baseline text\rule[-5pt]{2cm}{1ex}|: \\
  Baseline text\rule[-5pt]{2cm}{1ex}.\bigskip
\item \verb|\framebox[width][pos]{text}| \\
  e.g.: \verb|\framebox[9cm][s]{foo  bar baz}|:\\
  \framebox[9cm][s]{foo bar baz}\\\medskip
  where s=stretch.The other options are l=left and r=right
\end{itemize}
\end{frame}


\begin{frame}[fragile]\frametitle{Examples\dots}
\begin{itemize}
\item \verb|\makebox[width][pos]{text}| \\
  is the same as \verb|framebox|, but with no borders
\item \verb|\fbox{text}|\\
  is a short form: \verb|\fbox{Inspect this!}| \\
  \fbox{Inspect this!}
\item \verb|\mbox{text}|\\
  The invisible box \verb|\mbox{the contents}| \\ $\mapsto$
  \mbox{the contents} is treated as a single entity. This property is useful,
  and will be needed many times in the future\ldots
\end{itemize}
\end{frame}


\begin{frame}[fragile]\frametitle{Examples\dots}
\begin{itemize}
\item \verb|\vspace*{height} ja \hspace*{width}| \\
  \verb|\hspace*{0.3\textwidth}|\hspace*{0.3\textwidth} adds  
  certain amount of space to the documents\footnote{Observe the relative
  distance measure \texttt{\textbackslash textwidth}}\footnote{Footnotes can be
  produces with \texttt{\textbackslash footnote}\ldots}. Vertical space can be
  added with \verb|\vspace|. {\footnotesize\gray{(Negative distances
  move the following contents to the left or up. The command with star is 
  enforcing, otherwise adding the space
  is optional, and will be omitted at certain places, e.g. at the end of 
  paragraph)}}
\item \verb|\parbox[pos][height][ipos]{width}{text}| \\
  produces a box. For example, \\{\small
  \verb|The \fbox{\parbox[b][1cm][b]{4cm}{Schr�dinger's cat}}|
  The \fbox{\parbox[b][1cm][b]{4cm}{Schr�dinger's cat}}}
\end{itemize}
\end{frame}


\begin{frame}[fragile]\frametitle{Examples\dots}
\begin{itemize}
\item \verb|pos| can be t=top (line) or b=bottom (line)
\item \verb|ipos| can be t=top, b=bottom, c=centered or
  s=stretched
\item \verb|parbox| behaves like a small embedded
page\footnote{\texttt{\textbackslash minipage} will behave exactly like an
embedded page\ldots}, end it is treated as a single element.
\end{itemize}
\end{frame}


\begin{frame}[fragile]\frametitle{Examples\dots}
\begin{verbatim}
  \fbox{\parbox[t][3cm][t]{6cm}{%
      \includegraphics[width=1.6cm]
      {img/vttplain} \dotfill {\Huge $\Lambda$} \\
      \hspace*{1cm}\rotatebox{45}{%
        $\scriptstyle\begin{pmatrix} 
          a & b \\ c & d \end{pmatrix}$}}
  }
\end{verbatim}
\begin{center}
  \fbox{\parbox[t][3cm][t]{6cm}{%
      \includegraphics[width=1.6cm]
      {img/vttplain} \dotfill {\Huge $\Lambda$} \\
      \hspace*{1cm}\rotatebox{45}{%
        $\scriptstyle\begin{pmatrix} 
          a & b \\ c & d \end{pmatrix}$}}
  }
\end{center}
\end{frame}


\begin{frame}[fragile]\frametitle{Examples\dots}
What just happened?
\begin{itemize}
\item \verb|\fbox| included borders to the box \gray{(in contrast to mbox)}
\item \verb|\includegraphics| included the image \gray{(as pdf or png)}
\item \verb|\dotfill| fills the area with dots 
\item \verb|\hspace*| moved the <<pencil<< horizontally
\end{itemize}\bigskip

The following shows the same example with all elements having a box around 
them.
\end{frame}

\begin{frame}[fragile]\frametitle{Examples\dots}
Now we see what happens
\begin{center}
  \fbox{\parbox[t][3cm][t]{6cm}{%
      \fbox{\includegraphics[width=1cm]
        {img/vttplain}} \dotfill {\Huge \fbox{$\Lambda$}} \\
      \hspace*{1cm}\rotatebox{45}{\fbox{%
          $\scriptstyle\begin{pmatrix} 
            a & b \\ c & d \end{pmatrix}$}}
    }}
\end{center}\bigskip
The only difference is that the borders take some space, which makes a small
difference in the image.
\end{frame}


\begin{frame}[fragile]\frametitle{Conclusions}
We saw that \LaTeX{} code becomes complicated, if we
construct detailed elements. 

Because of this reason, the logical content of the document and fiddling with
the details should be kept \textbf{separated}.

In summary,
\begin{itemize}
\item everything can be tuned,
\item it is not wise to tune everything, 
\end{itemize}
\begin{flushleft}
but from experience I suppose that you would like to change the default font
sooner or later\ldots
\end{flushleft}
\end{frame}


\begin{frame}[fragile]\frametitle{Fonts}
The easiest way to get started with different document styles is to load a
pre-defined style package \gray{(from a \texttt{.sty} file)} in the manuscript
preamble:\smallskip
\begin{verbatim}
\usepackage{mystyle}
\end{verbatim}\smallskip
\begin{flushleft}
where the style is possibly one from the list \texttt{helvet, palatino,
avant, charter, bookman, newcent} or \texttt{times} \gray{(that are probably
already installed in your LaTeX-distribution package).}
\medskip

The size of the font can be changed with the descriptive commands
\verb|{\large  }| and \verb|{\small  }|.
\end{flushleft} 
\end{frame}


\begin{frame}[fragile]\frametitle{Fonts\dots}
\vspace*{-0.5cm}
\begin{tabular}{ll}
\verb|{\tiny Rabbit}| & {\tiny Rabbit} \\
\verb|{\scriptsize Rabbit}| & {\scriptsize Rabbit}  \\
\verb|{\footnotesize Rabbit}| & {\footnotesize Rabbit} \\ 
\verb|{\small Rabbit}| & {\small Rabbit}  \\ 
\verb|{\normalsize Rabbit}| & {\normalsize Rabbit}  \\
\verb|{\large Rabbit}| & {\large Rabbit} \\
\verb|{\Large Rabbit}| & {\Large Rabbit} \\
\verb|{\LARGE Rabbit}| & {\LARGE Rabbit} \\
\verb|{\huge Rabbit}| & {\huge Rabbit} \\
\verb|{\Huge Rabbit}| & {\Huge Rabbit} 
\end{tabular}
\end{frame}


\begin{frame}[fragile]\frametitle{Fonts\dots}
Observations
\begin{itemize}
\item These commands are declarations, not functions: \verb|\large| = change to
large font and keep in until the end of this (scope or) environment.
\item Anything between the \verb|\begin{}| \dots \verb|\end{}|
-block is an environment. A new environment can also be defined with the curly
braces \{ \dots \}.
\item Use the curly braces and think of the environments. This is better style
(from the programmer's point of view) than scattering the document with commands
like \verb|\normalsize| without a clear reference to the beginning of the
``non-normal'' behaviour.
\end{itemize}
\end{frame}


\begin{frame}[fragile]\frametitle{Fonts\dots}
I admit that even the last one \verb|\Huge| was not \textsl{that} large:\medskip
\begin{itemize}
\item[Q:] How to do a HUGE Rabbit?
\item[A:] \verb|\scalebox{4}{Rabbit}|
\end{itemize}\bigskip

\begin{center}
\scalebox{4}{Rabbit}
\end{center}
\end{frame}


\begin{frame}[fragile]\frametitle{Fonts\dots}
In addition, there are the ordinary functions to decorate the text\medskip

\begin{tabular}{ll}
\verb|\texttt{Rabbit}| & \texttt{Rabbit} \\
\verb|\textit{Rabbit}| & \textit{Rabbit} \\
\verb|\textbf{Rabbit}| & \textbf{Rabbit} \\
\verb|\textsl{Rabbit}| & \textsl{Rabbit} \\
\verb|\textsc{Rabbit}| & \textsc{Rabbit} \\
\verb|\underline{Rabbit}| & \underline{Rabbit} \\
\end{tabular}
\end{frame}


\begin{frame}[fragile]\frametitle{Fonts\dots}
These commands suffice well for a sophisticated writer.\medskip

More fonts can be found from the \LaTeX-distributions by searching the
\texttt{.fd} (font definition) files. For example,  
\texttt{t1pbk.fd} (T1 encoded postscript Bookman font definition)
can be selected with the definition
\begin{verbatim}
\fontfamily{pbk}\selectfont
This is Bookman Font
\end{verbatim}
{\fontfamily{pbk}\selectfont
This is bookman Font}
\end{frame}


\begin{frame}[fragile]\frametitle{Fonts\dots}
If you really need to change the font all the time \gray{(which is, in general,
a bad idea)}, it may be advantageous to define a custom command  in the
beginning of the document :\medskip
\begin{verbatim}
\newcommand{\avantgar}[1]{
{\fontfamily{pag}\selectfont #1}}
\end{verbatim}\medskip
etc\ldots
{\footnotesize
\begin{verbatim}
\newcommand{\bookman}[1]{{\fontfamily{pbk}\selectfont #1}}
\newcommand{\courier}[1]{{\fontfamily{pcr}\selectfont #1}}
\newcommand{\cmodern}[1]{{\fontfamily{cmr}\selectfont #1}}
\newcommand{\helvetic}[1]{{\fontfamily{phv}\selectfont #1}}
\newcommand{\newcent}[1]{{\fontfamily{pnc}\selectfont #1}}
\newcommand{\tmroman}[1]{{\fontfamily{ptm}\selectfont #1}}
\newcommand{\script}[1]{{\fontfamily{pzc}\selectfont #1}}
\end{verbatim}}
\end{frame}


\begin{frame}[fragile]\frametitle{Fonts\dots}
Now we have defined a couple of commands that fit the \LaTeX{}
philosophy.\medskip

\newcommand{\avantgar}[1]{{\fontfamily{pag}\selectfont #1}}
\newcommand{\bookman}[1]{{\fontfamily{pbk}\selectfont #1}}
\newcommand{\courier}[1]{{\fontfamily{pcr}\selectfont #1}}
\newcommand{\cmodern}[1]{{\fontfamily{cmr}\selectfont #1}}
\newcommand{\helvetic}[1]{{\fontfamily{phv}\selectfont #1}}
\newcommand{\newcent}[1]{{\fontfamily{pnc}\selectfont #1}}
\newcommand{\tmroman}[1]{{\fontfamily{ptm}\selectfont #1}}
\newcommand{\script}[1]{{\fontfamily{pzc}\selectfont #1}}

\begin{tabular}{ll}
\verb|\avantgar{Rabbit}| & \avantgar{Rabbit} \\
\verb|\bookman{Rabbit}| & \bookman{Rabbit} \\
\verb|\courier{Rabbit}| & \courier{Rabbit} \\
\verb|\cmodern{Rabbit}| & \cmodern{Rabbit} \\
\verb|\helvetic{Rabbit}| & \helvetic{Rabbit} \\
\verb|\tmroman{Rabbit}| & \tmroman{Rabbit} \\
%\verb|\utopia{Rabbit}| & \utopia{Rabbit} \\
\verb|\script{Rabbit}| & \script{Rabbit} \\
\end{tabular}\bigskip

\gray{\small
These commands are not included by default, because they will most likely ruin
your (technical or scientific) document. For the purposes you would be using
them, some of the office packages might be a better option.} 
\end{frame}


\begin{frame}[fragile]\frametitle{Conclusions\dots}
A good strategy for writing a paper with \LaTeX{} is to \medskip

\begin{itemize}
\item Make the title page by hand (e.g. by using the command just learned) 
\item Write the text with clean \LaTeX, using perhaps some of the extension
packages from AMS (American Mathematical Society) to stick with portability.
\item Make your own style sheet, if necessary
\item Include minor things in the document preamble before the
  \verb|\begin{document}| clause \gray{(even though it is sometimes legal to
  define command later in the document --- but what is the
  advantage?)}
\end{itemize}
\end{frame}


\begin{frame}[fragile]\frametitle{Conclusions\dots}
\begin{itemize}
\item \LaTeX{} and plain \TeX{} spaghetti mixed with the text
is a nightmare to maintain. \gray{Recall the experiences from HTML 3.2 before
properly working CSS.}
\item \LaTeX{} is a large software with a very long history (starting from
1978). Knowing the basics suffices well.
\end{itemize}
\end{frame}



\end{document}

